% ============================================================
% CAPÍTULO 2: FUNDAMENTOS ANTROPOLÓGICOS E ÉTICOS DA MH
% ============================================================

\chapter{Fundamentos Antropológicos e Éticos da \MH{}}
\label{cap:fundamentos}

\section{Introdução ao Capítulo}
\label{sec:intro-cap2}

A Encíclica \MH{} insere-se na tradição do magistério social da Igreja Católica, que desde a \textit{Rerum Novarum} (1891) vem oferecendo uma reflexão sistemática sobre as questões sociais, econômicas e políticas à luz do Evangelho e da razão natural. O que distingue a \MH{} neste contexto é seu objeto específico --- a dignidade e a vocação da pessoa humana na era da inteligência artificial --- e sua originalidade em articular uma antropologia integral que dialoga criticamente com os desafios tecnológicos contemporâneos.

Este capítulo examina os fundamentos antropológicos e éticos sobre os quais a Encíclica se assenta, com vistas a identificar os conceitos e princípios que poderão subsidiar a construção de um regime jurídico de proteção à personalidade humana. Para tanto, percorre-se o seguinte itinerário: a tradição do personalismo cristão como pano de fundo teórico (\ref{sec:personalismo}), a concepção de pessoa humana delineada pela \MH{} (\ref{sec:pessoa-mh}), a crítica ao paradigma tecnocrático (\ref{sec:paradigma-tecnocratico}) e os princípios éticos que orientam a relação entre técnica e humanidade (\ref{sec:principios-eticos}).

\section{A Tradição do Personalismo Cristão}
\label{sec:personalismo}

\subsection{Fundamentos Filosóficos do Personalismo}
\label{subsec:fundamentos-personalismo}

O personalismo constitui uma corrente filosófica que emergiu no século XX como reação tanto ao individualismo liberal quanto aos totalitarismos coletivistas. Seu núcleo central é a afirmação da pessoa como valor absoluto, irredutível a qualquer dimensão meramente funcional ou instrumental. Emmanuel Mounier, um de seus principais expoentes, define a pessoa como ``um ser espiritual, constituído como tal por uma forma de subsistência e de independência em seu ser; ela mantém essa subsistência mediante a adesão a uma hierarquia de valores livremente adotados, assimilados e vividos por um compromisso responsável'' \cite{Mounier1949P}.

O personalismo de Mounier distingue-se por recusar tanto o individualismo (que isola a pessoa em sua autonomia abstrata) quanto o coletivismo (que dissolve a pessoa em estruturas impessoais). A pessoa, nessa perspectiva, é essencialmente relacional: constitui-se na e pela relação com o outro, com a comunidade e com o transcendente. Essa dimensão relacional é o que distingue a pessoa do indivíduo: enquanto este é fechado em si mesmo, aquela se realiza no dom de si e na comunhão com os demais \cite{Ricoeur1990SM}.

Na tradição do pensamento cristão, essas intuições encontram lastro na teologia trinitária e na antropologia bíblica, segundo as quais o ser humano foi criado ``à imagem e semelhança de Deus'' (Gn 1,26-27), o que lhe confere uma dignidade intrínseca e inalienável. Tomás de Aquino, ao definir a pessoa como ``substância individual de natureza racional'', já estabelecia as bases para a compreensão da pessoa como um ser dotado de inteligência, vontade e capacidade de autodeterminação \cite{AquinoST}.

\subsection{O Personalismo no Magistério Social da Igreja}
\label{subsec:magisterio-social}

O Concílio Vaticano II representou um momento decisivo na incorporação do personalismo ao magistério católico. A Constituição Pastoral \textit{Gaudium et Spes} afirma que ``o homem, única criatura sobre a terra que Deus quis por si mesma, só pode encontrar-se plenamente a si mesmo pelo dom sincero de si mesmo'' \cite{ConcilioVaticanoII1965GS}. Essa formulação concilia a dignidade intrínseca da pessoa (``por si mesma'') com sua dimensão relacional (``dom sincero de si mesmo''), estabelecendo um paradigma que perpassa todo o magistério posterior.

João Paulo II, em sua encíclica \textit{Centesimus Annus}, aprofunda essa perspectiva ao afirmar que ``o homem vale mais pela sua \textit{essência} do que pela sua \textit{atuação}'' \cite{JoaoPauloII1991CA}. Essa distinção é crucial para o debate sobre a inteligência artificial, pois estabelece um critério antropológico fundamental: o valor da pessoa não deriva de sua utilidade, produtividade ou desempenho, mas de sua própria existência como ser dotado de dignidade.

Francisco, por sua vez, nas encíclicas \textit{Laudato Si'} e \textit{Fratelli Tutti}, estende essa reflexão às dimensões ecológica e social da vida contemporânea, denunciando as estruturas de poder que ameaçam a dignidade humana e propondo uma ecologia integral que articule justiça social, sustentabilidade ambiental e desenvolvimento humano \cite{Francisco2015LS, Francisco2020FT}.

\section{A Concepção de Pessoa Humana na \MH{}}
\label{sec:pessoa-mh}

\subsection{A Dignidade Ontológica da Pessoa}
\label{subsec:dignidade-ontologica}

A \MH{} retoma e desenvolve a tradição personalista do magistério, oferecendo uma concepção de pessoa humana articulada em três dimensões fundamentais: ontológica, relacional e transcendental.

A \textbf{dimensão ontológica} afirma que a pessoa humana possui um valor intrínseco que não depende de suas capacidades funcionais, de seu desempenho ou de sua utilidade social. A Encíclica rejeita expressamente toda forma de reducionismo que definiria o ser humano a partir de seus atributos mensuráveis ou de sua eficiência produtiva. Esse fundamento tem implicações diretas para o debate sobre a IA: um sistema algorítmico, por mais sofisticado que seja, não possui dignidade intrínseca, e sua operação deve estar sempre subordinada ao bem integral da pessoa \cite{LeaoXIV2026MH}.

A \textbf{dimensão relacional} da pessoa é igualmente central. A \MH{} recorda que o ser humano só pode realizar-se plenamente na relação com os outros e com o transcendente. Essa abertura constitutiva ao outro é o que distingue a pessoa de um sistema de IA, que opera exclusivamente com base em dados e algoritmos, sem capacidade genuína de saída de si em direção ao outro.

A \textbf{dimensão transcendental} reconhece na pessoa humana uma abertura ao infinito, uma sede de verdade e de bem que transcende as capacidades meramente computacionais. A Encíclica adverte contra a pretensão de reduzir a inteligência humana à inteligência artificial, como se esta fosse capaz de replicar ou superar a complexidade da consciência, da liberdade e da criatividade humanas.

\subsection{A Irredutibilidade da Pessoa aos Dados}
\label{subsec:irredutibilidade}

Um dos temas mais originais da \MH{} é a afirmação da irredutibilidade da pessoa humana aos seus dados. A Encíclica denuncia o que chama de ``reducionismo digital'': a tendência a conceber o ser humano como um mero conjunto de informações que podem ser coletadas, processadas e preditas por sistemas algorítmicos. Contra essa visão, o documento insiste que a pessoa é sempre mais do que seus dados, e que a subjetividade humana não pode ser integralmente traduzida em linguagem computacional.

Essa crítica dialoga diretamente com a literatura contemporânea sobre o capitalismo de vigilância. Shoshana Zuboff demonstrou como a extração e a análise de dados pessoais por grandes corporações tecnológicas constituem um novo modelo econômico que instrumentaliza a experiência humana como matéria-prima para predição e controle \cite{Zuboff2019SC}. A \MH{} oferece, no plano dos fundamentos, uma crítica antropológica a esse modelo: antes mesmo de ser ilegal, a redução da pessoa a seus dados é uma violação ontológica de sua dignidade.

\section{O Paradigma Tecnocrático}
\label{sec:paradigma-tecnocratico}

\subsection{Conceito e Crítica}
\label{subsec:conceito-paradigma}

A \MH{} dedica atenção especial ao que denomina ``paradigma tecnocrático'', conceito já presente na \textit{Laudato Si'} e que a nova Encíclica desenvolve e aprofunda. O paradigma tecnocrático é definido como a tendência a conceber a realidade exclusivamente sob o prisma da eficiência técnica, do cálculo utilitário e da manipulação instrumental, subordinando toda forma de saber e de valor à lógica da máquina.

A Encíclica não é ingênua quanto aos benefícios da tecnologia. Ela reconhece os avanços proporcionados pela inteligência artificial em áreas como medicina, educação e acesso à informação. A crítica dirigida não é à tecnologia em si, mas à mentalidade que absolutiza a técnica e a transforma em um fim em si mesmo, desvinculando-a de qualquer consideração ética ou antropológica.

A originalidade da \MH{} neste ponto reside em demonstrar que o paradigma tecnocrático não é apenas um problema ético, mas também antropológico e jurídico. Ao reduzir a pessoa a dados e as decisões a algoritmos, esse paradigma ataca a própria estrutura da personalidade humana, comprometendo a autonomia, a responsabilidade e a capacidade de deliberação moral que constituem o cerne da dignidade pessoal.

\subsection{Consequências do Paradigma Tecnocrático}
\label{subsec:consequencias-paradigma}

A \MH{} identifica três ordens de consequências decorrentes do paradigma tecnocrático que interessam diretamente à proteção jurídica da personalidade.

Em \textbf{primeiro lugar}, a \textbf{obscuridade decisória}. Sistemas de IA frequentemente operam como ``caixas pretas'', nas quais nem mesmo seus criadores conseguem explicar como determinada decisão foi tomada. Essa opacidade viola o direito fundamental à explicação das decisões automatizadas e compromete a possibilidade de controle e revisão que são inerentes ao Estado de Direito \cite{Pasquale2015BM, Wachter2017TSP}.

Em \textbf{segundo lugar}, a \textbf{desresponsabilização}. Quando decisões são delegadas a sistemas algorítmicos, cria-se uma ``difusão da responsabilidade'' que dificulta a identificação dos agentes responsáveis por eventuais danos. A \MH{} adverte que essa erosão da responsabilidade constitui uma das ameaças mais graves à pessoa humana, pois sem responsabilidade não há justiça, e sem justiça não há proteção efetiva da dignidade \cite{Jonas1979PR}.

Em \textbf{terceiro lugar}, a \textbf{homogeneização}. Sistemas algorítmicos tendem a classificar e categorizar as pessoas em perfis predeterminados, ignorando ou marginalizando aquilo que é singular, criativo ou imprevisível em cada ser humano. Essa tendência à homogeneização constitui uma agressão à unicidade da pessoa, que é um dos núcleos intangíveis dos direitos da personalidade.

\section{Os Princípios Éticos Fundamentais}
\label{sec:principios-eticos}

A \MH{}, articulando a tradição da doutrina social da Igreja com os desafios contemporâneos, propõe um conjunto de princípios éticos que devem orientar o desenvolvimento e a regulação da inteligência artificial. Estes princípios, embora formulados em linguagem teológica e filosófica, são suscetíveis de tradução em critérios normativos operacionais para o direito.

\subsection{O Princípio da Dignidade da Pessoa Humana}
\label{subsec:principio-dignidade}

O princípio fundamental que perpassa toda a \MH{} é o da dignidade da pessoa humana como valor absoluto e indisponível. A Encíclica reafirma que a pessoa humana ``não pode ser instrumentalizada por nenhum poder, seja ele econômico, político ou tecnológico''.

Esse princípio dialoga diretamente com o direito constitucional brasileiro, que elege a dignidade da pessoa humana como fundamento da República (art. 1º, III, da CF/88) e como princípio estruturante de todo o ordenamento jurídico \cite{Sarlet2021DF}. A \MH{} oferece, contudo, um aprofundamento do conteúdo desse princípio, explicitando suas implicações para o contexto tecnológico: a dignidade não é apenas um limite negativo (o que não se pode fazer com a pessoa), mas também um mandato positivo de promoção das condições para o pleno desenvolvimento da personalidade.

\subsection{O Princípio da Subsidiariedade}
\label{subsec:principio-subsidiariedade}

O princípio da subsidiariedade, clássico na doutrina social da Igreja, afirma que as instâncias superiores (Estado, organizações internacionais) não devem intervir naquilo que as instâncias inferiores (indivíduos, famílias, comunidades) podem realizar por si mesmas, mas têm o dever de apoiá-las quando necessário e de substituí-las quando insuficientes.

No contexto da inteligência artificial, a \MH{} aplica esse princípio para afirmar que a regulação da tecnologia não deve ser deixada exclusivamente à autorregulação das empresas ou às forças do mercado, mas requer uma ação coordenada dos poderes públicos que garanta a proteção da pessoa humana onde as instâncias privadas se mostram insuficientes ou inadequadas. Ao mesmo tempo, a subsidiariedade impede que o Estado assuma um controle excessivo sobre o desenvolvimento tecnológico, respeitando a liberdade de iniciativa e a autonomia dos corpos intermediários.

\subsection{O Princípio da Solidariedade}
\label{subsec:principio-solidariedade}

O princípio da solidariedade, na formulação da \MH{}, implica o reconhecimento de que o desenvolvimento tecnológico deve estar orientado para o bem comum e não apenas para o lucro ou o poder. A solidariedade exige que os benefícios da IA sejam distribuídos de forma equitativa e que os riscos não recaiam desproporcionalmente sobre os grupos mais vulneráveis.

Este princípio tem implicações jurídicas importantes: fundamenta a responsabilidade das empresas de tecnologia para com as comunidades afetadas por seus produtos, legitima políticas de inclusão digital e de acesso equitativo às tecnologias, e orienta a interpretação das normas de proteção de dados no sentido de proteger preferencialmente os grupos mais expostos a riscos algorítmicos \cite{Floridi2018AIG}.

\subsection{O Princípio da Transparência}
\label{subsec:principio-transparencia}

A transparência, para a \MH{}, não é apenas uma exigência técnica, mas um imperativo ético enraizado no reconhecimento da pessoa como sujeito racional que tem direito a compreender as decisões que a afetam. A opacidade algorítmica é denunciada como uma forma de dominação, pois ``o poder que não se deixa ver é um poder que escapa ao controle democrático''.

No plano jurídico, a transparência dos sistemas de IA se desdobra em três exigências normativas: (a) o direito à informação sobre o funcionamento dos sistemas algorítmicos; (b) o direito à explicação das decisões automatizadas individualizadas; e (c) o dever de auditoria independente dos sistemas de IA, especialmente daqueles classificados como de alto risco \cite{Wachter2017TSP}.

\subsection{O Princípio da Responsabilidade}
\label{subsec:principio-responsabilidade}

Finalmente, a \MH{} afirma o princípio da responsabilidade como correlato indispensável da liberdade e da criatividade humanas. Se o ser humano cria tecnologias poderosas, deve também responder por suas consequências. A Encíclica adverte contra a tentação de atribuir às máquinas uma autonomia que elas não possuem e de eximir os seres humanos da responsabilidade por suas criações.

O princípio da responsabilidade, na esteira de Hans Jonas \cite{Jonas1979PR}, adquire uma dimensão prospectiva: não se trata apenas de responder por danos já ocorridos (responsabilidade retrospectiva), mas de antecipar e prevenir os riscos antes que se concretizem (responsabilidade prospectiva). Essa dimensão é particularmente relevante para a regulação da IA, onde a incerteza sobre os efeitos futuros das tecnologias exige uma abordagem precaucional.

\section{Síntese do Capítulo}
\label{sec:sumario-cap2}

Os fundamentos antropológicos e éticos da \MH{}, conforme examinados neste capítulo, oferecem um referencial teórico consistente para a proteção jurídica da personalidade humana na era digital. A Encíclica, enraizada na tradição personalista e no magistério social da Igreja, propõe:

\begin{enumerate}
    \item Uma \textbf{antropologia integral} que afirma a dignidade ontológica, relacional e transcendental da pessoa humana, recusando todo reducionismo, inclusive o ``reducionismo digital'' que reduziria a pessoa a seus dados.
    \item Uma \textbf{crítica do paradigma tecnocrático} que identifica as ameaças à personalidade humana decorrentes da absolutização da técnica e da subordinação de toda realidade à lógica do cálculo e da eficiência.
    \item Um \textbf{conjunto de princípios éticos} --- dignidade, subsidiariedade, solidariedade, transparência e responsabilidade --- que podem ser traduzidos em critérios normativos para a regulação da IA.
\end{enumerate}

A hipótese que se coloca, e que será desenvolvida nos capítulos seguintes, é que esses fundamentos oferecem um referencial mais robusto para a proteção da personalidade do que as abordagens exclusivamente procedimentais ou economicistas que têm predominado no debate regulatório. Cabe agora examinar os desafios concretos que a IA impõe à personalidade humana (Capítulo 3), para depois avaliar como o direito brasileiro tem enfrentado essas questões (Capítulo 4) e, finalmente, propor as contribuições específicas da \MH{} para esse enfrentamento (Capítulo 5).
